\section{Statement of the problem and exact scope}

The two target pages mix four very different kinds of assertion.

\begin{enumerate}
\item There are exact finite identities, such as
\(G(n)=0!1!\cdots(n-2)!\) and the finite product of gamma values.
\item There are convergent infinite constructions, such as the canonical
product and the Taylor series.
\item There are analytic continuation statements involving a logarithm,
square root, quotient or path integral.  These are not meaningful until a
zero-free domain and a branch have been fixed.
\item There are asymptotic and uniqueness statements that quantify over a
precision, a sector, a differentiability order, or a coded function.
\end{enumerate}

A proof ``in iPRA'' must therefore say what finite arithmetical sentence replaces
each assertion.  It is not enough to copy a classical proof whose objects are
arbitrary Cauchy sequences or arbitrary functions.  Conversely, iPRA is not
limited to tiny finite computations.  A fixed approximation algorithm, an
explicit primitive-recursive modulus, and quantifier-free induction are enough
to express all of the finite estimates used below.

\subsection{Four status levels}\label{sec:scope-status}

Every entry in the coverage ledger is assigned one of the following levels.

\begin{description}
\item[\statusA: exact finite theorem.]  The claim reduces to natural-number,
integer or rational identities and inequalities.  The report gives a direct
primitive-recursive induction.  These are the first targets to compile into
\code{Prf PRAAx}.

\item[\statusB: quantitative analytic theorem.]  The report gives a finite
algorithm and an explicit error bound on a certified compact domain.  For every
fixed code and precision, the resulting statement is first-order arithmetic.
The remaining work is implementation of the checkers and their object-language
derivations.

\item[\statusC: branch- or representation-dependent theorem.]  The classical
formula is valid only after a continuation path, a branch of logarithm or
square root, or a normalization of a negative-order polygamma has been chosen.
The report supplies the missing data and proves the corresponding finite
contract.

\item[\statusD: corrected source statement.]  The displayed source formula is
false or misleading as written.  We identify the defect, prove the corrected
formula, and retain the original in the audit only.
\end{description}

None of these letters means ``already kernel-checked in Lean.''  A completed
object-iPRA theorem has the more specific form
\[
  []\;\vdash[\mathrm{PRAAx}]\;\varphi,
\]
where the terms in \(\varphi\) are built from the primitive-recursive symbols of
ProveIt's \code{ConciseIPRA.PR}.  The present package supplies the mathematics,
the exact target formulas and the dependency order, but does not claim that all
analytic targets have already been constructed as such derivations.

\subsection{Pinned repository foundation}

At the pinned ProveIt commit \cite{ProveIt}, iPRA is presented as first-order
intuitionistic logic with equality, a symbol for every primitive-recursive
function, its defining equations, successor axioms, and induction for
quantifier-free formulas.  The checked library has three relevant layers:

\begin{enumerate}
\item \code{ConciseIPRA} defines \code{PR}, terms, formulas, the axiom family
\code{PRAAx}, derivations \code{Prf}, and the bridge \code{pra\_iff\_native} to
the small native scalar calculus.
\item \code{IPRADerivations} supplies checked arithmetic: semiring laws,
order, bounded sums and products, division with remainder, pairing,
simultaneous recursion and polynomial normalization.
\item the \code{gamma\_ipra} packet gives a broad quantitative plan for Gamma,
including finite Euler products, explicit constants, integration, reflection,
multiplication and Stirling estimates.  Its own status statement correctly says
that the full analytic Lean development and the compilation into the object
theory remain work.
\end{enumerate}

This report reuses the third layer as a source of contracts, not as an assumed
theorem oracle.  Whenever a Barnes proof invokes Gamma, the formal dependency
is a named finite Gamma contract: an evaluator, a modulus, an apartness result,
or a fixed-sector Stirling remainder.  Such a contract can later be discharged
by the Gamma packet's intended modules.

\subsection{The theorem collection obtained}

The main conclusions are the following.

\begin{theorem}[Coverage theorem]\label{thm:coverage}
For every fixed primitive-recursive input name and every requested precision,
the following Barnes-$G$ statements admit an iPRA formulation by bounded
certificates and quantifier-free induction.

\begin{enumerate}[label=(\roman*)]
\item the canonical product, with a uniform compact-disc tail;
\item existence as an entire function, real conjugation, and the zero at
\(0,-1,-2,\ldots\) with multiplicities \(1,2,3,\ldots\);
\item normalization and the functional equation;
\item all positive-integer values, the superfactorial identity, the
hyperfactorial identity, and finite Gamma products;
\item the Taylor series for \(\log G(1+z)\), with explicit tails for every
closed subdisc of \(|z|<1\);
\item the logarithmic derivative, derivative, log-Gamma integral, and the
improper-integral representation on \(\Re z>-1\);
\item Kinkelin reflection, its Clausen forms, and the half-shift tangent form;
\item the squared absolute-value products, including the corrected
pure-imaginary formula;
\item the general multiplication formula, duplication, quarter values,
half-integer values, and the integer coefficient formula for all odd
half-integers;
\item the Hurwitz-zeta derivative representation, the negative-order
polygamma representation after normalization, and the analytic K-function
relation;
\item the full fixed-order Barnes asymptotic expansion on every fixed sector
separated from the negative real axis, with an explicit remainder bound;
\item order two and infinite type, with explicit integer-radius estimates;
\item the Vign\'eras characterization and the asymptotic characterization,
read as schemas for fixed coded functions with effective smoothness or
relative-error moduli.
\end{enumerate}
\end{theorem}

The proof occupies the rest of the report.  Its constructive content is not a
new axiom of real analysis.  It is the repeated pattern
\[
  \text{finite approximant}
  \quad+\quad
  \text{primitive-recursive error budget}
  \quad+\quad
  \text{bounded verifier}.
\]

\subsection{Two errors caught before formalization}

\begin{audit}[Asymptotic characterization]\label{audit:char}
The Wikipedia page presently displays
\(f(x+n)\sim\Gamma(x)^n n^{\binom{x}{2}}f(n)\).  Setting \(x=1\) contradicts
the functional equation, because the left side is exactly
\(\Gamma(n)f(n)\), while the displayed right side is \(f(n)\).  The cited
higher-order Bohr--Mollerup theorem has
\[
 f(x+n)\sim \Gamma(n)^x n^{\binom{x}{2}}f(n).
\]
We prove this corrected formula in \cref{sec:characterization}.
\end{audit}

\begin{audit}[Pure-imaginary modulus]\label{audit:modulus}
The Wikipedia page gives a factor \(y\) on the right of a formula for
\(|G(\ii y)|\) and says that \(y\) is arbitrary real.  The right side would be
negative for \(y<0\).  Conjugating the product gives a factor \(y^2\) in the
squared formula, and hence the unsquared factor is \(|y|\).  The corrected
identity is proved in \cref{sec:modulus}.
\end{audit}

These are not cosmetic matters.  The finite checker rejects both incorrect
statements immediately: the first at the exact input \(x=1\), the second at any
negative rational \(y\).  This illustrates why source auditing and formal
normalization should precede large proof-assistant developments.
